\documentclass[12pt,a4paper]{article}

% The report is designed for Tectonic's XeTeX engine and Overleaf.
\usepackage{fontspec}
\setmainfont{Liberation Serif}
\setmonofont{Liberation Mono}[Scale=MatchLowercase]
\usepackage[english]{babel}
\usepackage[left=2cm,right=1.5cm,top=2cm,bottom=2cm]{geometry}
\usepackage{setspace}
\usepackage{amsmath}
\usepackage{graphicx}
\usepackage{xcolor}
\usepackage{listings}
\usepackage{fancyvrb}
\usepackage{fvextra}
\usepackage{titlesec}
\usepackage{caption}
\usepackage{float}
\usepackage{array}
\usepackage{booktabs}
\usepackage[hidelinks]{hyperref}

\graphicspath{{figures/}}
\onehalfspacing
\setlength{\parindent}{1.25cm}
\setlength{\parskip}{0pt}
\setlength{\emergencystretch}{2em}

\definecolor{codebackground}{HTML}{F3F4F6}
\definecolor{codeframe}{HTML}{D1D5DB}
\definecolor{commentgreen}{HTML}{3F6212}
\definecolor{keywordblue}{HTML}{1D4ED8}

\titleformat{\section}
  {\normalfont\bfseries\centering\fontsize{13}{16}\selectfont}
  {\thesection}{0.8em}{\MakeUppercase}
\titleformat{name=\section,numberless}
  {\normalfont\bfseries\centering\fontsize{13}{16}\selectfont}
  {}{0pt}{\MakeUppercase}
\titleformat{\subsection}
  {\normalfont\bfseries\fontsize{12}{15}\selectfont}
  {\thesubsection}{0.7em}{}
\titlespacing*{\section}{0pt}{0pt}{12pt}
\titlespacing*{\subsection}{0pt}{10pt}{6pt}

\lstdefinestyle{wolfram}{
  language=Mathematica,
  basicstyle=\ttfamily\footnotesize,
  keywordstyle=\color{keywordblue}\bfseries,
  commentstyle=\color{commentgreen},
  morecomment=[s]{(*}{*)},
  morekeywords={RandomPrime,RandomInteger,PowerMod,IntegerLength,IntegerDigits,
    IntegerString,FromDigits,ToCharacterCode,FromCharacterCode,StringJoin,GCD,
    PrimeQ,Hash,ByteArray,SymmetricKey,Encrypt,Decrypt,Mod,While,If,Break},
  backgroundcolor=\color{codebackground},
  frame=single,
  rulecolor=\color{codeframe},
  breaklines=true,
  columns=fullflexible,
  keepspaces=true,
  showstringspaces=false,
  numbers=left,
  numberstyle=\tiny\color{gray},
  xleftmargin=0.5cm,
  framexleftmargin=0.4cm,
  aboveskip=7pt,
  belowskip=7pt
}

% \screenshot{file name without .png}{caption}
% Shows the screenshot if it is in figures/, otherwise a placeholder with the expected name.
\newcommand{\screenshot}[2]{%
  \begin{figure}[H]
    \centering
    \IfFileExists{figures/#1.png}{%
      \includegraphics[width=0.92\textwidth,height=0.55\textheight,keepaspectratio]{#1.png}%
    }{%
      \fbox{\parbox[c][3cm][c]{0.85\textwidth}{\centering
        \texttt{figures/\detokenize{#1}.png}\\[4pt]\textit{screenshot not added yet}}}%
    }
    \caption{#2}
    \label{fig:#1}
  \end{figure}%
}

\begin{document}

\begin{titlepage}
  \begin{center}
    \IfFileExists{figures/LOGO_UTM.jpg}{%
      \includegraphics[width=3.4cm]{figures/LOGO_UTM.jpg}\par
      \vspace{0.25cm}
    }{}

    {\bfseries\fontsize{14}{17}\selectfont
      MINISTRY OF EDUCATION AND RESEARCH\par
      OF THE REPUBLIC OF MOLDOVA\par
      Technical University of Moldova\par
      Faculty of Computers, Informatics and Microelectronics\par
      Department of Software Engineering and Automation\par
    }

    \vspace{0.8cm}
    {\bfseries\fontsize{16}{19}\selectfont Daniel Moisei, FAF-241\par}

    \vspace{0.8cm}
    {\bfseries\fontsize{36}{40}\selectfont Report\par}
    \vspace{0.25cm}
    {\itshape\fontsize{20}{24}\selectfont Laboratory Work No. 5\par}
    {\bfseries\itshape\fontsize{20}{24}\selectfont
      \href{https://else.fcim.utm.md/course/section.php?id=11371}{Cryptography and Security}\par}

    \vspace{0.7cm}
    {\fontsize{20}{25}\selectfont Topic: Public-Key Cryptography\par}
    \vspace{0.2cm}
    {\fontsize{15}{19}\selectfont RSA, ElGamal and the Diffie--Hellman Key Exchange\par}
  \end{center}

  \vfill
  \begin{flushright}
    {\fontsize{15}{18}\selectfont Checked by:}\par
    \textbf{Maia Zaica}, \textit{university assistant}
  \end{flushright}

  \vfill
  \begin{center}
    {\bfseries\fontsize{14}{17}\selectfont Chi\c{s}in\u{a}u -- 2026}
  \end{center}
\end{titlepage}

\pagenumbering{arabic}
\setcounter{page}{2}
\tableofcontents
\newpage

\section*{Introduction}
\addcontentsline{toc}{section}{Introduction}

The purpose of this laboratory work is to study public-key (asymmetric) cryptography and to apply three of its classical algorithms to real-size numbers. The tasks are:

\begin{itemize}
  \item \textbf{Task 2.1} -- generate RSA keys with a modulus $n$ of at least 2048 bits, then encrypt and decrypt the message $m = $ \texttt{Moisei Daniel};
  \item \textbf{Task 2.2} -- generate ElGamal keys for the 2048-bit prime $p$ and the generator $g = 2$ given in the assignment, then encrypt and decrypt the same message;
  \item \textbf{Task 3} -- perform the Diffie--Hellman key exchange between Alisa and Bob over the same $p$ and $g$, with randomly chosen secrets $a$ and $b$, and obtain a 256-bit key for AES.
\end{itemize}

All computations were done in Wolfram Language (Wolfram Cloud notebook), which works with arbitrarily large integers exactly. Every step was checked inside the notebook: the bit sizes of the keys, the key equations, and that decryption returns the original message. The screenshots of the notebook are included in each section; the complete values (600-digit numbers) are listed in Appendix~A.

\newpage
\section{Short Theory}

\subsection{Public-Key Cryptography}

In symmetric cryptography both parties share one secret key, which has to be delivered over a secure channel. Public-key cryptography solves this problem with a pair of keys: a \emph{public key}, which can be published, and a \emph{private key}, which only the owner knows. Anyone can encrypt with the public key, but only the owner of the private key can decrypt.

The security of these systems relies on \emph{one-way functions}: operations that are easy to compute but practically impossible to invert without extra information. Two such problems are used here:
\begin{itemize}
  \item \textbf{integer factorization} -- multiplying two large primes is easy, but recovering them from their product is infeasible (RSA);
  \item \textbf{discrete logarithm} -- computing $g^x \bmod p$ is easy, but finding $x$ from $g^x \bmod p$ is infeasible for a large prime $p$ (ElGamal, Diffie--Hellman).
\end{itemize}

Both rely on fast modular exponentiation $a^b \bmod n$ (square-and-multiply), which in Wolfram Language is \texttt{PowerMod[a, b, n]}. With a negative exponent, \texttt{PowerMod[a, -1, n]} returns the modular inverse $a^{-1} \bmod n$, computed with the extended Euclidean algorithm.

\subsection{RSA}

\textbf{Key generation.} Choose two different large primes $p$ and $q$ and compute
\begin{equation}
  n = pq, \qquad \varphi(n) = (p-1)(q-1).
\end{equation}
Choose a public exponent $e$ with $1 < e < \varphi(n)$ and $\gcd(e, \varphi(n)) = 1$, and compute the private exponent
\begin{equation}
  d \equiv e^{-1} \pmod{\varphi(n)}, \qquad \text{i.e. } ed \equiv 1 \pmod{\varphi(n)}.
\end{equation}
The public key is $(n, e)$; the private key is $(n, d)$; $p$, $q$ and $\varphi(n)$ are kept secret.

\textbf{Encryption and decryption.} For a message $0 \le m < n$:
\begin{equation}
  c = m^e \bmod n, \qquad m = c^d \bmod n.
\end{equation}
Decryption works because $ed = 1 + k\varphi(n)$, so by Euler's theorem $c^d = m^{ed} = m \cdot (m^{\varphi(n)})^k \equiv m \pmod n$.

\subsection{ElGamal}

ElGamal works in the multiplicative group modulo a large prime $p$ with a generator $g$.

\textbf{Key generation.} Choose a random private key $x$ and compute the public key $y = g^x \bmod p$. The public key is $(p, g, y)$, the private key is $x$.

\textbf{Encryption.} For a message $0 < m < p$ choose a fresh random ephemeral key $k$ and compute
\begin{equation}
  c_1 = g^k \bmod p, \qquad c_2 = m \cdot y^k \bmod p.
\end{equation}
The ciphertext is the pair $(c_1, c_2)$. The value $s = y^k = g^{xk}$ acts as a one-time mask on the message.

\textbf{Decryption.} The receiver recomputes the mask with the private key and removes it:
\begin{equation}
  s = c_1^x \bmod p = g^{kx} \bmod p, \qquad m = c_2 \cdot s^{-1} \bmod p.
\end{equation}
The key $k$ must be new for every message: if it is reused, $c_2/c_2' = m/m'$ and one known message reveals the other.

\subsection{Diffie--Hellman Key Exchange}

The Diffie--Hellman protocol lets two parties agree on a common secret over a public channel. Both know $p$ and $g$.
\begin{enumerate}
  \item Alisa chooses a random secret $a$ and sends $A = g^a \bmod p$ to Bob.
  \item Bob chooses a random secret $b$ and sends $B = g^b \bmod p$ to Alisa.
  \item Alisa computes $K = B^a \bmod p$, Bob computes $K = A^b \bmod p$.
\end{enumerate}
Both obtain the same value, because $B^a = g^{ba} = g^{ab} = A^b \pmod p$. An eavesdropper sees only $p$, $g$, $A$ and $B$; to obtain $K$ it would have to solve a discrete logarithm. Plain Diffie--Hellman does not authenticate the parties, so in practice the public values are signed to prevent a man-in-the-middle attack.

\subsection{From the Shared Secret to an AES-256 Key}

The shared secret $K$ is a 2048-bit number, while AES-256 needs a key of exactly 256 bits. The secret is therefore passed through a key derivation step: $K$ is written as a fixed-length sequence of 256 bytes (big-endian, keeping leading zeros) and hashed with SHA-256, whose output is exactly 256 bits:
\begin{equation}
  k_{\mathrm{AES}} = \mathrm{SHA\text{-}256}\bigl(\mathrm{bytes}_{256}(K)\bigr).
\end{equation}
Since Alisa and Bob have the same $K$, they derive the same AES key.

\subsection{Group Parameters from the Assignment}

For Tasks 2.2 and 3 the assignment gives a 2048-bit prime $p$ (617 decimal digits, listed in Appendix~A) and $g = 2$. This is the standard \texttt{ffdhe2048} group from RFC~7919. The prime is a \emph{safe prime}: $q = (p-1)/2$ is also prime. Because $2^q \equiv 1 \pmod p$, the generator $g = 2$ produces the subgroup of prime order $q$, so the secret exponents ($x$, $k$, $a$, $b$) are chosen at random from $\{2, \ldots, q-1\}$.

\newpage
\section{Message Encoding}

As required by the assignment, the message $m = $ \texttt{Moisei Daniel} (surname and name, separated by one space) is converted to a number in three steps: each character is replaced by its ASCII code, each code is written as two hexadecimal digits, and the resulting hexadecimal string is read as one integer and converted to decimal.

\begin{table}[H]
  \centering
  \caption{ASCII encoding of the message}
  \renewcommand{\arraystretch}{1.2}
  \setlength{\tabcolsep}{4.2pt}
  \begin{tabular}{l*{13}{c}}
    \toprule
    Character & M & o & i & s & e & i & \textit{sp} & D & a & n & i & e & l \\
    ASCII     & 77 & 111 & 105 & 115 & 101 & 105 & 32 & 68 & 97 & 110 & 105 & 101 & 108 \\
    Hex       & 4D & 6F & 69 & 73 & 65 & 69 & 20 & 44 & 61 & 6E & 69 & 65 & 6C \\
    \bottomrule
  \end{tabular}
\end{table}

\begin{align}
  H &= \texttt{4D6F697365692044616E69656C}_{16}, \\
  m &= 6135048831842424644607035401580_{10}.
\end{align}

The number $m$ has 31 decimal digits (103 bits), so it is smaller than the RSA modulus and the ElGamal prime, as both algorithms require. After decryption the inverse conversion is applied: the integer is split into base-256 digits (bytes) and each byte is turned back into its character.

\begin{lstlisting}[style=wolfram,caption={Message encoding (cell 1)}]
plainText  = "Moisei Daniel";
asciiCodes = ToCharacterCode[plainText, "ASCII"];                 (* character -> ASCII code *)
hexBytes   = ToUpperCase[IntegerString[#, 16, 2]] & /@ asciiCodes; (* code -> two hex digits *)
messageHex = StringJoin[hexBytes];                                 (* bytes in original order *)
m          = FromDigits[messageHex, 16];                           (* hex string -> integer m *)
decode[x_Integer] := FromCharacterCode[IntegerDigits[x, 256], "ASCII"];
\end{lstlisting}

\screenshot{fig01_message_encoding}{Message encoding: ASCII codes, hex bytes, the integer $m$ and the round-trip check}

\newpage
\section{Task 2.1. RSA}

\subsection{Key Generation}

\begin{enumerate}
  \item The public exponent is fixed to $e = 65537 = 2^{16}+1$, the usual choice: it is prime and has only two bits set, so encryption is fast.
  \item Two random primes $p$ and $q$ are taken from $[2^{1023}, 2^{1024}-1]$ with \texttt{RandomPrime}, so each has exactly 1024 bits and their product has 2047 or 2048 bits.
  \item $n = pq$ and $\varphi(n) = (p-1)(q-1)$ are computed. The generation is repeated until $p \ne q$, $n$ has exactly 2048 bits and $\gcd(e, \varphi(n)) = 1$ (otherwise $e$ would have no inverse).
  \item The private exponent is $d = e^{-1} \bmod \varphi(n)$, computed with \texttt{PowerMod[e, -1, phi]}.
\end{enumerate}

\begin{lstlisting}[style=wolfram,caption={RSA key generation (cell 2)}]
rsaE = 65537;
While[True,
  rsaP   = RandomPrime[{2^1023, 2^1024 - 1}];
  rsaQ   = RandomPrime[{2^1023, 2^1024 - 1}];
  rsaN   = rsaP*rsaQ;                         (* n = p q *)
  rsaPhi = (rsaP - 1)*(rsaQ - 1);             (* phi(n) *)
  If[rsaP != rsaQ && IntegerLength[rsaN, 2] == 2048 && GCD[rsaE, rsaPhi] == 1, Break[]]
];
rsaD = PowerMod[rsaE, -1, rsaPhi];            (* d = e^-1 mod phi(n) *)
\end{lstlisting}

The output (Figure~\ref{fig:fig02_rsa_keygen}) shows the first and last digits of each number and its size: $p$ and $q$ have 1024 bits, $n$ has 2048 bits, $\gcd(e,\varphi(n)) = 1$ and $ed \bmod \varphi(n) = 1$, which confirms that $d$ is the correct inverse.

\screenshot{fig02_rsa_keygen}{RSA key generation and the checks of the keys}

\subsection{Encryption and Decryption}

\begin{enumerate}
  \item The condition $0 \le m < n$ is checked, since RSA works with residues modulo $n$.
  \item Encryption with the public key: $c = m^e \bmod n$.
  \item Decryption with the private key: $m' = c^d \bmod n$.
  \item $m'$ is compared with $m$ and converted back to text.
\end{enumerate}

\begin{lstlisting}[style=wolfram,caption={RSA encryption and decryption (cell 3)}]
rsaCipher    = PowerMod[m, rsaE, rsaN];          (* c  = m^e mod n *)
rsaRecovered = PowerMod[rsaCipher, rsaD, rsaN];  (* m' = c^d mod n *)
\end{lstlisting}

\screenshot{fig03_rsa_encrypt_decrypt}{RSA encryption, decryption and recovery of the message}

The ciphertext $c$ is a 2048-bit number that bears no visible relation to $m$. Decryption returned exactly $m = 6135048831842424644607035401580$, which decodes to \texttt{Moisei Daniel}.

\newpage
\section{Task 2.2. ElGamal}

\subsection{Group Parameters}

The prime $p$ from the assignment is entered in the notebook and $g = 2$, $q = (p-1)/2$ are defined. The checks confirm that $p$ has 2048 bits, that $p$ and $q$ are prime, and that $g^q \bmod p = 1$, i.e.\ $g$ generates the subgroup of order $q$.

\begin{lstlisting}[style=wolfram,caption={Group parameters (cell 4, $p$ shortened)}]
p = 3231700607131100730015351347782516...4281003994966956119696956248629032338072839127039;
g = 2;
q = (p - 1)/2;                                (* p = 2q + 1 is a safe prime *)
\end{lstlisting}

\screenshot{fig04_group_parameters}{The prime $p$ and generator $g$ from the assignment, with their checks}

\subsection{Key Generation}

The private key $x$ is a random integer from $\{2,\ldots,q-1\}$ and the public key is $y = g^x \bmod p$.

\begin{lstlisting}[style=wolfram,caption={ElGamal key generation (cell 5)}]
elgX = RandomInteger[{2, q - 1}];              (* private key x *)
elgY = PowerMod[g, elgX, p];                   (* public key y = g^x mod p *)
\end{lstlisting}

\screenshot{fig05_elgamal_keygen}{ElGamal public key $(p, g, y)$ and private key $x$}

\subsection{Encryption and Decryption}

\begin{enumerate}
  \item A fresh ephemeral key $k$ is chosen at random from $\{2,\ldots,q-1\}$.
  \item The first part of the ciphertext is $c_1 = g^k \bmod p$; it lets the receiver rebuild the mask.
  \item The mask is $s = y^k \bmod p$ and the second part of the ciphertext is $c_2 = m \cdot s \bmod p$.
  \item The receiver computes the same mask from its private key: $s' = c_1^x \bmod p = g^{kx} \bmod p$.
  \item The message is recovered with the modular inverse of the mask: $m' = c_2 \cdot s'^{-1} \bmod p$. The inverse exists because $p$ is prime and $s' \ne 0$.
\end{enumerate}

\begin{lstlisting}[style=wolfram,caption={ElGamal encryption and decryption (cell 6)}]
elgK  = RandomInteger[{2, q - 1}];              (* ephemeral key k *)
elgC1 = PowerMod[g, elgK, p];                   (* c1 = g^k mod p *)
elgS  = PowerMod[elgY, elgK, p];                (* mask s = y^k mod p *)
elgC2 = Mod[m*elgS, p];                         (* c2 = m s mod p *)

elgS2        = PowerMod[elgC1, elgX, p];        (* s' = c1^x mod p *)
elgRecovered = Mod[elgC2*PowerMod[elgS2, -1, p], p];  (* m' = c2 s'^-1 mod p *)
\end{lstlisting}

\screenshot{fig06_elgamal_encrypt_decrypt}{ElGamal ciphertext $(c_1, c_2)$, the mask check and the recovered message}

The masks computed by the sender and by the receiver are equal, and the recovered integer equals $m$ and decodes to \texttt{Moisei Daniel}. Unlike RSA, the ElGamal ciphertext is twice as long as the modulus and is different each time, even for the same message, because $k$ is random.

\newpage
\section{Task 3. Diffie--Hellman Key Exchange}

\subsection{Exchange between Alisa and Bob}

Alisa and Bob use the public parameters $p$ and $g = 2$ from the assignment.

\begin{enumerate}
  \item Alisa chooses a random secret $a \in \{2,\ldots,q-1\}$ and Bob chooses $b \in \{2,\ldots,q-1\}$ independently.
  \item They compute and exchange the public values $A = g^a \bmod p$ and $B = g^b \bmod p$.
  \item Alisa computes $K = B^a \bmod p$ and Bob computes $K = A^b \bmod p$.
  \item The two results are compared: they must be equal to $g^{ab} \bmod p$.
\end{enumerate}

\begin{lstlisting}[style=wolfram,caption={Diffie--Hellman exchange (cell 7)}]
aliceSecret = RandomInteger[{2, q - 1}];            (* a *)
bobSecret   = RandomInteger[{2, q - 1}];            (* b *)
alicePublic = PowerMod[g, aliceSecret, p];          (* A = g^a mod p, sent to Bob *)
bobPublic   = PowerMod[g, bobSecret, p];            (* B = g^b mod p, sent to Alisa *)

aliceShared = PowerMod[bobPublic, aliceSecret, p];  (* Alisa: K = B^a mod p *)
bobShared   = PowerMod[alicePublic, bobSecret, p];  (* Bob:   K = A^b mod p *)
\end{lstlisting}

\screenshot{fig07_dh_exchange}{Secrets, public values and the shared secret computed by Alisa and by Bob}

\subsection{Deriving the AES-256 Key}

Each party writes its shared secret as 256 bytes (\texttt{IntegerDigits[K, 256, 256]}, which keeps leading zero bytes) and hashes it with SHA-256. The 32-byte digest is the AES-256 key.

\begin{lstlisting}[style=wolfram,caption={AES key derivation (cell 8)}]
toBytes[x_Integer] := ByteArray[IntegerDigits[x, 256, 256]];   (* 256-byte big-endian *)
aliceKey  = Hash[toBytes[aliceShared], "SHA256", "ByteArray"]; (* 32 bytes = 256 bits *)
bobKey    = Hash[toBytes[bobShared],   "SHA256", "ByteArray"];
aesKeyHex = Hash[toBytes[aliceShared], "SHA256", "HexString"];
\end{lstlisting}

\screenshot{fig08_dh_aes_key}{The 256-bit AES key derived by Alisa and Bob}

Both parties obtained the same 64-hex-digit (256-bit) key:
\begin{center}
  \texttt{2f8b1b3890e49f6c55b1f322fb1887a5aa461fd6beba991dd47db6f75317d740}
\end{center}

\subsection{Using the Key with AES}

To show that the agreed key is usable, Alisa encrypts the message with AES-256 under her key and Bob decrypts it with the key he derived on his side.

\begin{lstlisting}[style=wolfram,caption={AES-256 encryption with the agreed key (cell 9)}]
aliceAES  = SymmetricKey[<|"Cipher" -> "AES256", "Key" -> aliceKey|>];
bobAES    = SymmetricKey[<|"Cipher" -> "AES256", "Key" -> bobKey|>];
aesCipher = Encrypt[aliceAES, plainText];
Decrypt[bobAES, aesCipher]
\end{lstlisting}

\screenshot{fig09_aes_demo}{AES-256 encryption by Alisa and decryption by Bob with the agreed key}

\newpage
\section{Results}

\begin{table}[H]
  \centering
  \caption{Summary of the checks performed in the notebook}
  \renewcommand{\arraystretch}{1.25}
  \begin{tabular}{>{\raggedright\arraybackslash}p{0.17\textwidth}>{\raggedright\arraybackslash}p{0.47\textwidth}>{\raggedright\arraybackslash}p{0.25\textwidth}}
    \toprule
    \textbf{Task} & \textbf{Check} & \textbf{Result} \\
    \midrule
    Encoding & \texttt{decode[m] == "Moisei Daniel"} & True \\
    2.1 RSA & bit length of $n$ & 2048 \\
            & $\gcd(e, \varphi(n))$, \ $ed \bmod \varphi(n)$ & 1, \ 1 \\
            & $c^d \bmod n = m$ & True, text recovered \\
    Parameters & bits of $p$; $p$, $q$ prime; $g^q \bmod p$ & 2048; True, True; 1 \\
    2.2 ElGamal & sender mask $=$ receiver mask & True \\
                & $c_2 s^{-1} \bmod p = m$ & True, text recovered \\
    3 DH & $B^a \bmod p = A^b \bmod p$ & True \\
         & SHA-256 keys of Alisa and Bob equal & True, 256 bits \\
         & AES-256 decryption by Bob & \texttt{Moisei Daniel} \\
    \bottomrule
  \end{tabular}
\end{table}

\section{Conclusions}

This laboratory work applied three public-key algorithms to numbers of real cryptographic size. The message \texttt{Moisei Daniel} was converted through its ASCII and hexadecimal representation into the integer $m = 6135048831842424644607035401580$, which was then encrypted and decrypted with RSA and with ElGamal.

For RSA a 2048-bit modulus was built from two random 1024-bit primes, and the private exponent was obtained as the modular inverse of $e = 65537$. The checks $\gcd(e,\varphi(n)) = 1$ and $ed \equiv 1 \pmod{\varphi(n)}$ confirmed the keys, and decryption returned the original message. The security of RSA rests on the fact that $d$ can only be computed from $\varphi(n)$, which requires the factorization of $n$.

For ElGamal and Diffie--Hellman the 2048-bit safe prime from the assignment was used; the notebook confirmed that $p$ and $(p-1)/2$ are prime and that $g = 2$ generates the subgroup of order $q$. ElGamal hides the message with a random mask $g^{xk}$ that only the owner of $x$ can rebuild, so the same message gives a different ciphertext each time. In the Diffie--Hellman exchange Alisa and Bob published only $A$ and $B$, yet both computed the same 2048-bit secret, from which SHA-256 produced the same 256-bit AES key; an AES-256 encryption by Alisa was correctly decrypted by Bob.

In practice public-key algorithms are much slower than symmetric ones, so they are combined: RSA or Diffie--Hellman is used to establish a key, and the data itself is encrypted with a symmetric cipher such as AES -- exactly the scheme demonstrated in Task 3.

\section{Source Code}

The complete Wolfram Language code, with one commented block per notebook cell, is available in the course repository:

\begin{center}
  \url{https://github.com/Moisei-D/CS-Labs/tree/main/labs/Lab05}
\end{center}

\newpage
\appendix
\renewcommand{\thesection}{\Alph{section}}
\section{Complete Numerical Values}

The values below were exported from the notebook. Long integers are wrapped over several lines.

\IfFileExists{wolfram_output.txt}{%
  \VerbatimInput[fontsize=\scriptsize,breaklines=true,breakanywhere=true,frame=single,rulecolor=\color{codeframe}]{wolfram_output.txt}%
}{%
  \begin{center}\fbox{\parbox{0.85\textwidth}{\centering
    \texttt{wolfram\_output.txt}\\[4pt]\textit{paste the output of cell 10 into this file}}}\end{center}%
}

\end{document}
